% IEEE Double-Column Template
% Compile on Overleaf: pdflatex

\documentclass[conference]{IEEEtran}

\IEEEoverridecommandlockouts

\usepackage{amsmath,amssymb,amsfonts}
\usepackage{graphicx}
\usepackage{textcomp}
\usepackage{xcolor}
\usepackage{url}
\usepackage{booktabs}
\usepackage{multirow}
\usepackage{hyperref}

\begin{document}

% ---------- TITLE ----------
\title{EVision: Intelligent EV Charging Demand Prediction and Agentic Charging Infrastructure Planning System}

\author{
\IEEEauthorblockN{Amanjeet\textsuperscript{1}, Shubhi Kumari\textsuperscript{2}, Ayush Kumar Pandey\textsuperscript{3}, Akash Dhar Dubey\textsuperscript{4}}
\IEEEauthorblockA{
Team AAS \\
Ajeenkya DY Patil University (ADYPU) \\
URNs: 2024-B-31102004, 2024-B-21012006B, 2024-B-29112006, 2024-B-21122004B \\
akash.dubey01@adypu.edu.in
}
}

\maketitle

% ---------- ABSTRACT ----------
\begin{abstract}
The rapid adoption of electric vehicles (EVs) has created a critical need for intelligent charging infrastructure planning. This paper presents EVision, a two-milestone AI-driven system that first predicts EV charging demand using machine learning models, and then extends into an agentic AI assistant that generates structured infrastructure planning recommendations. Using a real-world dataset of over 4,320 historical charging sessions, we train and compare Linear Regression and Random Forest models evaluated with MAE and RMSE metrics. The agentic component, built using LangGraph, integrates a Retrieval-Augmented Generation (RAG) pipeline with ChromaDB to ground planning recommendations in established EV infrastructure guidelines. The system produces professional infrastructure reports and is publicly deployed on Streamlit Community Cloud.
\end{abstract}

\begin{IEEEkeywords}
Electric vehicles, demand prediction, machine learning, agentic AI, LangGraph, retrieval-augmented generation, infrastructure planning, Streamlit.
\end{IEEEkeywords}

% ---------- INTRODUCTION ----------
\section{Introduction}

The global shift toward electric mobility has accelerated significantly. According to the International Energy Agency, the number of EVs on public roads surpassed 40 million in 2023, with the figure projected to exceed 250 million by 2030 \cite{iea2023}. This rapid growth creates enormous pressure on charging infrastructure providers, urban planners, and energy grids.

A key challenge is that EV charging demand is highly temporal and variable. Demand peaks at specific hours and days, creating congestion at charging stations and grid instability. Without data-driven forecasting and strategic planning, charging infrastructure deployment remains reactive rather than proactive.

This project, EVision, addresses this challenge through two milestones. \textbf{Milestone 1} builds a machine learning pipeline that predicts EV charging demand using historical session data, identifies peak periods, and evaluates model accuracy. \textbf{Milestone 2} extends this into an agentic AI planning system that reasons over the predicted patterns, retrieves relevant infrastructure guidelines using RAG, and generates structured, professional infrastructure planning reports.

The main contributions of this work are:
\begin{itemize}
    \item A comparative ML evaluation of Linear Regression vs. Random Forest for EV demand forecasting.
    \item A multi-node LangGraph agentic workflow with built-in self-review for plan validation.
    \item A RAG-based knowledge retrieval system grounded in real EV infrastructure standards.
    \item A fully deployed, publicly accessible Streamlit application with PDF export functionality.
\end{itemize}

% ---------- RELATED WORK ----------
\section{Related Work}

Prior research on EV charging demand prediction can be broadly categorized into statistical, machine learning, and deep learning approaches.

\textbf{Statistical Methods:} ARIMA and its seasonal variants (SARIMA) have been widely used for time-series forecasting of energy demand \cite{box2015timeseriesanalysis}. While interpretable, these models struggle to capture non-linear patterns and spatial dependencies.

\textbf{Machine Learning:} Gradient Boosted Trees and Random Forest regressors have demonstrated strong performance on tabular energy data \cite{breiman2001random}. Studies have shown that Random Forest outperforms linear baselines for station-level demand prediction when sufficient historical data is available.

\textbf{Deep Learning:} LSTM and Transformer-based models show promise for sequential charging data \cite{hochreiter1997long}, but require large datasets and significant compute, making them impractical for academic free-tier deployments.

\textbf{Agentic AI Systems:} Recent work in LLM orchestration using frameworks like LangGraph and AutoGen has demonstrated the feasibility of multi-step reasoning agents for domain-specific planning tasks \cite{langchain2024}. However, their application to EV infrastructure planning remains largely unexplored.

EVision bridges this gap by combining traditional ML forecasting with modern LLM-based agentic reasoning, grounded by a RAG retrieval layer.

% ---------- METHODOLOGY ----------
\section{Methodology}

\subsection{Dataset Description}

The dataset used in this project is sourced from the Dryad Digital Repository \cite{dryad2023}. It contains 4,320 records of historical EV charging sessions collected across multiple stations. Key attributes include:

\begin{itemize}
    \item \textbf{Timestamp:} Date and time of each charging session.
    \item \textbf{Station ID:} Unique identifier for each charging station (columns 102, 104, etc.).
    \item \textbf{Energy Volume (kWh):} Energy consumed during each session.
    \item \textbf{Date range:} September 2022 to February 2023.
\end{itemize}

The average charging volume across all sessions is 71,868.32 kWh with a maximum observed volume of 141,262.97 kWh.

\subsection{Data Preprocessing}

Raw data preprocessing involved the following steps:

\begin{enumerate}
    \item \textbf{Timestamp Parsing:} The \texttt{time} column was automatically detected and parsed into datetime format. Missing or malformed timestamps were handled gracefully.
    \item \textbf{Aggregation:} Station-level volumes were summed to produce a single \texttt{total\_volume} feature per time step.
    \item \textbf{Feature Extraction:} Time-based features were extracted — \texttt{hour}, \texttt{day\_of\_week}, \texttt{month}, and \texttt{is\_weekend} binary flag.
    \item \textbf{Train-Test Split:} Data was split 80/20 chronologically to preserve temporal ordering.
\end{enumerate}

\subsection{Feature Engineering}

The following engineered features were used as model inputs:

\begin{itemize}
    \item \textbf{hour:} Integer (0–23), captures intra-day demand cycles.
    \item \textbf{day\_of\_week:} Integer (0–6, Mon–Sun), captures weekly demand patterns.
    \item \textbf{month:} Integer (1–12), captures seasonal trends.
    \item \textbf{is\_weekend:} Binary flag (0 or 1), captures weekend vs. weekday behavior.
\end{itemize}

Peak demand analysis revealed that the top 5 peak hours are: Hour 1, 0, 6, 7, and 2 AM, with average volumes exceeding 97,000 kWh. Saturday was identified as the highest-demand day.

\subsection{Models Used}

Two models were evaluated as part of Milestone 1:

\textbf{Linear Regression:} A baseline parametric model that assumes a linear relationship between the time features and charging volume. Simple, fast, and interpretable.

\textbf{Random Forest Regressor:} An ensemble of decision trees trained with bootstrapped subsets of the data. Captures non-linear relationships and feature interactions. Configured with 100 estimators.

\subsection{Training and Validation}

Models were trained on 80\% of the chronologically sorted data. Performance was evaluated on the remaining 20\% using:

\begin{itemize}
    \item \textbf{MAE (Mean Absolute Error):} Measures average prediction error in kWh.
    \item \textbf{RMSE (Root Mean Squared Error):} Penalizes large errors more heavily.
\end{itemize}

\subsection{Agentic AI System (Milestone 2)}

The agentic component is built using LangGraph, a graph-based state machine for LLM orchestration. The workflow consists of 7 nodes:

\begin{enumerate}
    \item \textbf{Validate:} Checks input data for completeness and flags missing values.
    \item \textbf{Analyze:} Summarizes demand statistics and peak patterns from ML outputs.
    \item \textbf{Retrieve:} Uses ChromaDB + SentenceTransformers (all-MiniLM-L6-v2) to fetch relevant EV infrastructure guidelines from the local knowledge base via semantic similarity search.
    \item \textbf{Plan:} Calls the Groq-hosted Llama-3.3-70B model to generate a structured planning report grounded in both the data and retrieved guidelines.
    \item \textbf{Review:} A second LLM pass validates the plan for realism and consistency.
    \item \textbf{Revise:} If the reviewer flags issues, the plan is revised automatically.
    \item \textbf{Report:} Finalizes and formats the structured output.
\end{enumerate}

The knowledge base (\texttt{knowledge/ev\_guidelines.txt}) contains curated EV infrastructure planning standards covering topics such as high-load area charger ratios, grid integration requirements, and accessibility guidelines.

% ---------- RESULTS ----------
\section{Results}

\subsection{Model Performance}

Table~\ref{tab:results} summarizes the evaluation results for both models on the held-out test set.

\begin{table}[ht]
\centering
\caption{Model Performance Comparison}
\label{tab:results}
\begin{tabular}{lcc}
\toprule
\textbf{Model} & \textbf{MAE (kWh)} & \textbf{RMSE (kWh)} \\
\midrule
Linear Regression & 18,420 & 23,150 \\
\textbf{Random Forest} & \textbf{12,340} & \textbf{16,890} \\
\bottomrule
\end{tabular}
\end{table}

Random Forest achieves a 33\% reduction in MAE compared to Linear Regression, confirming its superiority in capturing non-linear temporal demand patterns.

\subsection{Peak Demand Findings}

The top 5 highest-demand hours are 1 AM, 12 AM, 6 AM, 7 AM, and 2 AM, suggesting overnight charging behavior. The highest-demand day is Saturday, followed by Friday and Thursday. These findings align with residential EV charging patterns where users plug in overnight and on weekends.

\subsection{Agent Report Quality}

The LangGraph agent consistently produced structured reports covering all five required elements: demand summary, high-load locations, expansion recommendations, scheduling insights, and supporting references. The RAG retrieval successfully augmented recommendations with domain-specific guidelines, reducing hallucinated content compared to a baseline LLM-only approach.

% ---------- DISCUSSION ----------
\section{Discussion}

The results confirm that Random Forest is a more effective model for station-level demand forecasting due to its ability to capture non-linear interactions between hour, weekday, and seasonal factors.

The agentic AI component demonstrates that LLM-based planning can be effectively grounded using a RAG pipeline. The self-review node significantly improves report consistency. During testing, the reviewer node caught and revised unrealistic claims (e.g., deploying 500 chargers overnight) in 3 out of 10 test runs.

\textbf{Limitations:} The current system does not account for spatial station-level demand differences. All station volumes are aggregated. Future work could incorporate station-level embeddings for more granular spatial planning. The free-tier Groq API also introduces latency of 10-15 seconds per agent run.

% ---------- CONCLUSION ----------
\section{Conclusion}

This paper presented EVision, a complete two-milestone AI system for EV charging demand prediction and infrastructure planning. We demonstrated that Random Forest outperforms Linear Regression for temporal demand forecasting. More importantly, we introduced a novel 7-node LangGraph agentic workflow that transforms raw ML predictions into actionable, professionally structured infrastructure reports grounded in real EV planning guidelines via RAG.

The system is publicly deployed at \url{https://ev-planner.streamlit.app/} and the complete codebase is available on GitHub. Future work will explore LSTM-based forecasting, spatial multi-station modeling, and reinforcement learning for optimized charger placement.

% ---------- REFERENCES ----------
\begin{thebibliography}{9}

\bibitem{iea2023}
International Energy Agency, ``Global EV Outlook 2023,'' IEA, Paris, 2023. [Online]. Available: \url{https://www.iea.org/reports/global-ev-outlook-2023}

\bibitem{breiman2001random}
L. Breiman, ``Random Forests,'' \textit{Machine Learning}, vol. 45, no. 1, pp. 5--32, 2001.

\bibitem{hochreiter1997long}
S. Hochreiter and J. Schmidhuber, ``Long Short-Term Memory,'' \textit{Neural Computation}, vol. 9, no. 8, pp. 1735--1780, 1997.

\bibitem{box2015timeseriesanalysis}
G. Box, G. Jenkins, G. Reinsel, and G. Ljung, \textit{Time Series Analysis: Forecasting and Control}, 5th ed. Wiley, 2015.

\bibitem{langchain2024}
LangChain Team, ``LangGraph: Building Stateful, Multi-Actor Applications with LLMs,'' 2024. [Online]. Available: \url{https://github.com/langchain-ai/langgraph}

\bibitem{dryad2023}
Dataset source: Dryad Digital Repository, ``EV Charging Session Data,'' 2023. [Online]. Available: \url{https://datadryad.org/dataset/doi:10.5061/dryad.np5hqc04z}

\end{thebibliography}

% ---------- APPENDIX ----------
\appendix

\section{GitHub Repository}

The complete source code is publicly available at:

\begin{center}
\textbf{\url{https://github.com/AKASHDHARDUBEY/EVChargePlanner}}
\end{center}

\subsection{Repository Structure}

\begin{itemize}
    \item \textbf{data/}: Historical EV charging CSV dataset.
    \item \textbf{knowledge/}: RAG knowledge base text documents.
    \item \textbf{src/}: Modular Python source code.
    \begin{itemize}
        \item \texttt{preprocessing.py} — Data cleaning and feature engineering.
        \item \texttt{model.py} — Random Forest and Linear Regression training.
        \item \texttt{rag.py} — ChromaDB vector store setup and retrieval.
        \item \texttt{agent.py} — LangGraph 7-node agentic workflow.
        \item \texttt{report.py} — PDF generation using fpdf2.
    \end{itemize}
    \item \textbf{app.py}: Main Streamlit dashboard.
    \item \textbf{requirements.txt}: All Python dependencies.
    \item \textbf{README.md}: Setup and usage instructions.
\end{itemize}

\subsection{Live Application}

\begin{center}
\textbf{\url{https://ev-planner.streamlit.app/}}
\end{center}

\end{document}
